% Generated by scripts/assemble.py; corrections belong in source/corrections.json.
\examhead{35}{复旦大学 2026 年分析试卷}{38}
% p039-b04; source PDF page 39
\blockanchor{p039-b04}
\question{1}已知函数项级数 $\sum_{n=2}^{\infty} \frac{\cos(n\pi x)}{n^{ax}}$ 关于 $x \in (2, +\infty)$ 一致收敛, 求实数 $a$ 的取值范围.\par
% p039-b05; source PDF page 39
\blockanchor{p039-b05}
\question{2}计算积分 $\int_0^{+\infty}\left(\frac{x^2}{1 + x^4}\right)^2\mathrm{d}x.$\par
% p039-b06; source PDF page 39
\blockanchor{p039-b06}
\question{3}求曲面 $z = xy$ 被柱面 $x^{2} + y^{2} = 1$ 所截部分的面积.\par
% p039-b07; source PDF page 39
\par\Needspace{6\baselineskip}
\blockanchor{p039-b07}
\question{4}已知椭圆 $C = \{(x, y) \mid x^2 + 4y^2 = 16\}$, 方向取逆时针, 计算曲线积分\par
\nopagebreak[4]
% p039-b08; source PDF page 39
\blockanchor{p039-b08}
\displayformula{\int_ {C} \left(\frac {x y ^ {2}}{x ^ {2} + 2 y ^ {2} + 1} - y\right) \mathrm{d} x + \left(x + \frac {x ^ {2} y}{x ^ {2} + 2 y ^ {2} + 1}\right) \mathrm{d} y.}
% p039-b09; source PDF page 39
\blockanchor{p039-b09}
\question{5}求 $\sin \ln x$ 在 $x = 1$ 处带 Peano 余项的四阶 Taylor 展开式.\par
% p039-b10; source PDF page 39
\blockanchor{p039-b10}
\question{6}已知 $f(x) \in C[0,1]$ 在 (0,1) 上二阶可导, $f(0) = f(1) = 0$, $f\left(\frac{1}{2}\right) = \frac{1}{4}$, 证明: 存在 $\xi \in (0,1)$, 使得 $f''(\xi) = -2$.\par
% p039-b11; source PDF page 39
\blockanchor{p039-b11}
\question{7}已知 $\zeta(x) = \sum_{n=1}^{\infty} \frac{1}{n^x}$, 证明: $\zeta(x)$ 在 $(1, +\infty)$ 上任意阶导数存在, 并判断 $\int_{1}^{+\infty} (\zeta(x) - 1) \, \mathrm{d}x$ 的敛散性.\par
% p039-b12; source PDF page 39
\par\Needspace{6\baselineskip}
\blockanchor{p039-b12}
\question{8}计算极限\par
\nopagebreak[4]
% p039-b13; source PDF page 39
\blockanchor{p039-b13}
\displayformula{\lim _ {\varepsilon \rightarrow 0 ^ {+}} \iint_ {x ^ {2} + x y + y ^ {2} \geq \varepsilon^ {2}} \frac {1}{x ^ {2} + y ^ {2}} \sin \frac {1}{x ^ {2} + y ^ {2}} \mathrm{d} x \mathrm{d} y.}
% p039-b14; source PDF page 39
\blockanchor{p039-b14}
并说明计算过程的合理性.\par
% p039-b15; source PDF page 39
\blockanchor{p039-b15}
\question{9}证明幂级数 $\sum_{n=0}^{\infty} \frac{x^{3n}}{(3n)!}$ 的和函数是微分方程 $\frac{\mathrm{d}^2 y}{\mathrm{d}x^2} + \frac{\mathrm{d}y}{\mathrm{d}x} + y = e^x$ 的一个解, 并求出和函数的具体表达式.\par
% p039-b16; source PDF page 39
\blockanchor{p039-b16}
\question{10}证明: $(0,0)$ 是自洽系统 $\left\{ \begin{array}{l} \frac{\mathrm{d}x}{\mathrm{d}t} = y - f(x) \\ \frac{\mathrm{d}y}{\mathrm{d}t} = -g(x) \end{array} \right.$ 的奇点, 并判断其李雅普诺夫意义下的稳定性, 其中 $f, g$ 是实值连续可微函数, 满足 $x \neq 0$ 时, $xf(x) > 0$, $xg(x) > 0$.\par
% p039-b17; source PDF page 39
\blockanchor{p039-b17}
\question{11}设 $f(z)$ 在区域 $\Omega$ 上全纯, 满足 $\sin (f(z)) = f(z)$, 证明: $f(z)$ 是常值函数.\par
% p039-b18; source PDF page 39
\blockanchor{p039-b18}
\question{12}设 $f(z)$ 在复平面上全纯, 且对任意 $\operatorname{Re} z \neq 0$ 的 $z \in \mathbb{C}$ 都有 $|f(z)| < \frac{1}{|\operatorname{Re} z|}$. 证明: $f(z) = 0$.\par
% p039-b19; source PDF page 39
\blockanchor{p039-b19}
\question{13}设 $f, g$ 是实直线 $\mathbb{R}$ 上的勒贝格可积函数, 证明: 二元函数 $F(x, y) = f(x)g(y + x)$ 是 $\mathbb{R}^2$ 上的勒贝格可积函数.\par
% p039-b20; source PDF page 39
\par\Needspace{6\baselineskip}
\blockanchor{p039-b20}
\question{14}证明黎曼勒贝格引理: 对于 $[a, b]$ 上的勒贝格可积函数 $f$ 以及 $[0, 1]$ 上的勒贝格可积的周期函数 $g$, 有\par
\nopagebreak[4]
% p039-b21; source PDF page 39
\blockanchor{p039-b21}
\displayformula{\lim _ {n \rightarrow \infty} \int_ {a} ^ {b} f (x) g (n x) \mathrm{d} x = \int_ {0} ^ {1} g (x) \int_ {a} ^ {b} f (x) \mathrm{d} x.}
